\DocumentMetadata{
  pdfversion=2.0,
  pdfstandard=ua-2,
  lang=en-US,
  tagging=on,
  tagging-setup={math/setup=mathml-SE,tabsorder=structure}
}
\documentclass[11pt,letterpaper]{article}
\usepackage[margin=0.68in,includefoot]{geometry}
\usepackage{newcomputermodern}
\usepackage{amsmath,amscd,mathtools}
\providecommand{\square}{\mdlgwhtsquare}
\usepackage[hidelinks]{hyperref}
\hypersetup{
  pdftitle={Analysis First-Year Exam, August 2018, Part 2},
  pdfauthor={University of Florida Department of Mathematics},
  pdfsubject={Graduate mathematics examination},
  pdfdisplaydoctitle=true,
  bookmarksopen=true
}
\setcounter{secnumdepth}{0}
\setlength{\parindent}{0pt}
\setlength{\parskip}{0.28em}
\setlength{\emergencystretch}{3em}
\setlength{\leftmargini}{2.15em}
\setlength{\leftmarginii}{2.65em}
\setlength{\leftmarginiii}{3em}
\renewcommand{\labelenumi}{\textbf{\theenumi.}}
\pagestyle{empty}
\raggedbottom
\allowdisplaybreaks
\newenvironment{examproblems}[1]{%
  \begin{enumerate}%
  \setcounter{enumi}{#1}%
  \setlength{\topsep}{0.35em}%
  \setlength{\partopsep}{0pt}%
  \setlength{\itemsep}{0.48em}%
  \setlength{\parsep}{0.18em}%
}{\end{enumerate}}
\newenvironment{examsubparts}{%
  \begin{enumerate}%
  \setlength{\topsep}{0.18em}%
  \setlength{\partopsep}{0pt}%
  \setlength{\itemsep}{0.16em}%
  \setlength{\parsep}{0.1em}%
}{\end{enumerate}}
\newenvironment{examinstructions}{%
  \par\begingroup\itshape
}{%
  \par\endgroup\vspace{0.18em}
}
\makeatletter
\renewcommand\section{\@startsection{section}{1}{0pt}%
  {-1.25ex plus -.25ex minus -.1ex}%
  {0.55ex plus .1ex}%
  {\normalfont\large\bfseries}}
\renewcommand{\@maketitle}{%
  \newpage\null\vskip -1.2em%
  {\footnotesize\bfseries
    \href{https://www.ufl.edu/}{University of Florida}, \href{https://math.ufl.edu/}{Department of Mathematics}\par}%
  \vskip 0.18em%
  \tagstructbegin{tag=Title}%
    {\LARGE\bfseries\href{https://gma.math.ufl.edu/exams/analysis-first-year/}{Analysis First-Year Exam}, August 2018, Part 2\par}%
  \tagstructend%
  \vskip 0.35em%
  \tagmcbegin{artifact}%
    \hrule height 1pt%
  \tagmcend%
  \vskip 0.55em%
}
\makeatother
\title{Analysis First-Year Exam, August 2018, Part 2}
\author{}
\date{}
\begin{document}
\maketitle
\thispagestyle{empty}
\begin{examinstructions}
Answer FOUR questions in detail. State carefully any results used without proof.
\end{examinstructions}
\begin{examproblems}{0}
\item
Let the sequence \((s_n)_{n=1}^{\infty}\) in \([0,1]\) be uniformly distributed in the sense that if \(0 \le a \le b \le 1\) then 
\[
\lim_{n\to\infty}\frac{\#\{k\le n:s_k\in[a,b]\}}{n}=b-a.
\]
 Let \(f:[0,1]\to\mathbb{R}\) and prove that 
\[
\lim_{n\to\infty}\frac{f(s_1)+\cdots+f(s_n)}{n}=\int_0^1 f(t)\,dt
\]
 in each of the following cases:
\begin{examsubparts}
\item[\textnormal{(i)}]
\(f\) is a step function (a finite linear combination of indicators of intervals);
\item[\textnormal{(ii)}]
\(f\) is Riemann-integrable.
\end{examsubparts}
\item
Fix \(a\in\mathbb{R}\) and for each integer \(n>0\) write \(f_n(t)=n^a t(1-t^2)^n\) whenever \(0\le t\le 1\).
\begin{examsubparts}
\item[\textnormal{(i)}]
Show that \((f_n)_{n=1}^{\infty}\) converges pointwise on \([0,1]\); say to~\(f\).
\item[\textnormal{(ii)}]
For which values of~\(a\) does \((f_n)_{n=1}^{\infty}\) converge uniformly on \([0,1]\)? Justify.
\item[\textnormal{(iii)}]
For which values of~\(a\) is it true that \(\int_0^1 f_n\to\int_0^1 f\)? Justify.
\end{examsubparts}
\item
Let \(f:[1,\infty)\to\mathbb{R}\) be continuous and satisfy \(\lim_{t\to\infty}f(t)=A\in\mathbb{R}\). Prove that there is a sequence \((p_n)_{n=0}^{\infty}\) of polynomials such that \(p_n(1/t)\) converges to \(f(t)\) uniformly for \(t\ge 1\).
\item
Let \((f_n)_{n=1}^{\infty}\) be a sequence of real-valued measurable functions. Prove that each of the following sets is measurable:
\begin{examsubparts}
\item[\textnormal{(i)}]
\(B=\{\omega:(f_n(\omega))_{n=1}^{\infty}\text{ has no biggest term}\}\);
\item[\textnormal{(ii)}]
\(C=\{\omega:\cos(f_n(\omega))>0\text{ for each }n>0\}\);
\item[\textnormal{(iii)}]
\(D=\{\omega:(f_n(\omega))_{n=1}^{\infty}\text{ does not converge to a rational number}\}\).
\end{examsubparts}
\item
State the Monotone Convergence Theorem. Use it to prove the following: let \((f_n)_{n=1}^{\infty}\) be a sequence of non-negative integrable functions with pointwise limit~\(f\); if \(\int f_n\,d\mu\le M<\infty\) for each~\(n\) then~\(f\) is integrable and \(\int f\,d\mu\le M\).
\end{examproblems}
\end{document}
